\documentclass[english,twoside, openright, hidelinks, 11 pt, class=report,crop=false]{standalone}
\input{../../preamb}
\input{../bm}
\begin{document}

\opgt

\op{opgsumtotall}
Skriv tallene som summen av to tall. \os
\opgeks{
3 kan skrives som $ 1+2 $
}\vspace{5pt}
\abch{
\item 4
\item 5 
\item 6
\item 7
\item 8
\item 9
}

\op{opgsumtretall}
Skriv tallene som summen av tre tall. \os
\opgeks{
	4 kan skrives som $ 1+2+1 $
} \vspace{5pt}
\abch{
	\item 5 
	\item 6
	\item 7
	\item 8
	\item 9
}

\op{opgtiervenn}
To heltall som til sammen utgjør 10 kalles \outl{tiervenner}. For eksempel er 1 og 9 tiervenner fordi $ 1+9=10 $. \os

\abcn{
\item Finn tiervennen til \os
\abch{
	\item 2
	\item 3 
	\item 4
	\item 5
} \vsk

\item  Når man har gjort oppgave 1), hvorfor er det ikke ''nødvendig'' å finne tiervennene til 6, 7 og 8? 
}
\newpage
\op{opgsumparodd}
\mers{Du kan tillate deg å svare på spørsmålene bare ved å prøve ut et par\\ eksempler. For bevis, se \refgrub{opgaddifparodd}.}\os

Velg rett alternativ av 1), 2) og 3).
\abc{
	\item Summen av to partall er \vspace{-5pt}
	\abcn{
		\item et partall.\\
		\item et oddetall.\\
		\item noen ganger et partall, andre ganger et oddetall.
	}
	\item Summen av to oddetall er \vspace{-5pt}
	\abcn{
		\item et partall.\\
		\item et oddetall.\\
		\item noen ganger et partall, andre ganger et oddetall.
	}
	\item Summen av et partall og et oddetall er \vspace{-5pt}
	\abcn{
		\item et partall.\\
		\item et oddetall.\\
		\item noen ganger et partall, andre ganger et oddetall.
	}
}

\nes
\op{opgdifto}
Skriv tallene som differansen mellom to tall.
\opgeks{1 kan skrives som $ 8-7 $.} \vspace{5pt}
\abch{
	\item 2
	\item 3
	\item 4
	\item 5
	\item 6
	\item 7
	\item 8
}
\newpage
\op{opgaddifomv}
Fyll ut det manglende svaret ved å utnytte at addisjon og subtraksjon er omvendte operasjoner.
\opgeks{
Siden $9+7=16$, er $16-7=9$ og $16-9=7$. \os
Siden $13-8=5$, er $5+8=13$ og $13-8=5$.
}
\abc{
\item Siden $8+13=21$, er $21-13=....$ og $21-8=....$
\item Siden $42+19=61$, er $61-19=....$ og $61-42=....$
\item Siden $762+141=903$, er $903-141=....$ og $903-762=....$
\item Siden $15-9=6$, er $6+9=....$ og $15-6=....$.
\item Siden $83-25=58$, er $58+25=....$ og $83-58=....$.
\item Siden $904-352=552$, er $552+352=....$ og $904-552=....$
}

\op{opgdifparodd}
{\footnotesize (I denne oppgaven kan du tillate deg å svare på spørsmålene bare ved å prøve ut et par eksempler. For bevis, se oppgave \refops{opgaddifparodd}.)}
\abc{
	\item Differansen mellom to partall er \vspace{-5pt}
	\abcn{
		\item Et partall.\\
		\item Et oddetall.\\
		\item Noen ganger et partall, andre ganger et oddetall.
	}
	\item Differansen mellom to oddetall er \vspace{-5pt}
	\abcn{
		\item Et partall.\\
		\item Et oddetall.\\
		\item Noen ganger et partall, andre ganger et oddetall.
	}
	\item Differansen mellom et partall og et oddetall er \vspace{-5pt}
	\abcn{
		\item Et partall.\\
		\item Et oddetall.\\
		\item Noen ganger et partall, andre ganger et oddetall.
	}
}
\newpage
\nes
\op{opgadsomgang}
Skriv som gangestykker og alternativ sum. \os
\opgeks{$ 3+3+3+3+3=3\cdot 5=5+5+5 $}
\abc{
	\item $ 2+2+2 $
	\item $ 3+3+3+3+3+3 $
	\item $ 4+4 $ \\
	\item $ 5+5+5+5+5+5+5+5+5+5 $
	\item $ 6+6+6+6 $
	\item $ 7+7+7+7 $
}

\op{opggangrut}
Tegn ruter og finn svaret på gangestykket.
\opgeks{\vsb
	\begin{center}
		\includegraphics{\figp{opggangruteks1}}\hspace{3cm}
		\includegraphics{\figp{opggangruteks2}}
	\end{center}
} \vspace{5pt}

\abch{
	\item $ 4\cdot 5 $ 
	\item $ 3\cdot 8 $
	\item $ 2\cdot 9 $
	\item $ 5\cdot 6 $
	\item $ 7\cdot 8$
}
\newpage
\op{opggangdivomv}
Fyll ut de manglende svarene ved å utnytte at
multiplikasjon og divisjon er omvendte operasjoner.
\opgeks{
Siden $3\cdot 8=24$, er $24:8=3$ og $24:3=8$. \os
Siden $12:3=4$, er $3\cdot4=12$ og $12:4=3$.
}
\abc{
\item Siden $9\cdot 5=45$, er $45:9=....$ og $45:5=....$
\item Siden $23\cdot8=184$, er $184:23=....$ og $184:8=....$
\item Siden $42\cdot56=2352$, er $2352:42=....$ og $2352:=....$
\item Siden $21:3=7$, er $3\cdot7=....$ og $21:7=....$
\item Siden $216:9=24$, er $9\cdot24=....$ og $216:24=....$
\item Siden $3796:52=73$, er $52\cdot73=....$ og $3796:73=....$
}

\newpage
\grubop{opggangparod} \vs
\abc{
	\item Vil et heltall ganget med 2 alltid resultere i et partall eller et oddetall?
	\item Vil et partall ganget med 5 alltid resultere i et partall eller et oddetall? Hvilket siffer vil alltid stå på enerplassen?
	\item Vil et oddetall ganget med 5 alltid resultere i et partall eller et oddetall? Hvilket siffer vil alltid stå på enerplassen?
}


\end{document}